\documentclass[12pt,a4paper]{article}
\usepackage{amsmath,amssymb,amsthm}
\usepackage{hyperref}
\usepackage{geometry}
\usepackage{enumitem}
\usepackage{booktabs}
\geometry{margin=2.5cm}

\newtheorem{theorem}{Theorem}[section]
\newtheorem{corollary}[theorem]{Corollary}
\newtheorem{proposition}[theorem]{Proposition}
\theoremstyle{definition}
\newtheorem{definition}[theorem]{Definition}
\theoremstyle{remark}
\newtheorem{remark}[theorem]{Remark}

\title{The Top Quark Mass from Recognition Science:\\
Rung 21 on the $\varphi$-Ladder}

\author{Jonathan Washburn\\
Recognition Science Research Institute\\
Austin, Texas\\
\texttt{washburn.jonathan@gmail.com}}

\date{March 2026}

\begin{document}
\maketitle

\begin{abstract}
We derive the top quark mass from the Recognition Science
$\varphi$-ladder.  The top quark is the heaviest fermion, sitting at
up-type quark rung $r_t = 21$ (base rung $4$ plus second-generation
spacing $\tau(2) = 17$).  The master mass law gives
$m_t = \mathrm{yardstick}(\mathrm{UpQuark}) \cdot \varphi^{21 - 8
+ \mathrm{gap}(Q_t)}$, where $Q_t = +2/3$ is the top quark charge.
The yardstick for the up-type sector is
$2^{-1} E_{\mathrm{coh}} \varphi^{35}$, with $B_{\mathrm{pow}} = -1$
and $r_0 = 35$.  The experimental value
$m_t = 172.69~\mathrm{GeV}$~\cite{PDG} is recovered through the
calibration seam.  We explain why the top Yukawa coupling $y_t
\approx 1$ is natural in RS (the top sits near the electroweak scale
on the $\varphi$-ladder) and derive the top-bottom mass ratio
$m_t/m_b \approx \varphi^8$.  All structural results are derived from the RS
axioms~\cite{RS_Axioms}.
\end{abstract}

%% ============================================================
\section{Introduction}
\label{sec:intro}
%% ============================================================

The top quark was discovered in 1995 by the CDF and D0 collaborations
at the Fermilab Tevatron~\cite{CDF1995, D01995}.  With a mass of
$m_t = 172.69 \pm 0.30~\mathrm{GeV}$~\cite{PDG}, it is the heaviest
fundamental fermion---nearly as massive as a gold atom.  In the
Standard Model, the top mass is set by the Yukawa coupling $y_t
\approx 1$ to the Higgs field: $m_t = y_t v/\sqrt{2}$, where
$v \approx 246~\mathrm{GeV}$ is the electroweak vacuum expectation
value.  Unlike the other quarks, whose Yukawa couplings span six
orders of magnitude ($y_u \sim 10^{-5}$, $y_c \sim 10^{-2}$, $y_t
\sim 1$), the top coupling is ``natural'' in the sense that it is
order unity.

Recognition Science (RS)~\cite{RS_Axioms} derives the top quark
mass from the $\varphi$-ladder with no free Yukawa coupling.  The golden ratio
$\varphi = (1+\sqrt{5})/2$ and coherence quantum $E_{\mathrm{coh}} =
\varphi^{-5}~\mathrm{eV}$ emerge from the Recognition Composition
Law; the top occupies a definite rung $r_t = 21$ on the ladder, fixed
by $D = 3$ cube geometry.  The top's special status---its proximity
to the electroweak scale---is structural, not accidental.

%% ============================================================
\section{Recognition Science Framework}
\label{sec:framework}
%% ============================================================

Recognition Science derives all physics from a single primitive,
the \emph{recognition cost functional} $J(x)$, uniquely determined
by three axioms.

\begin{definition}[RS Cost Functional]
For $x > 0$: $J(x) = \tfrac{1}{2}(x + x^{-1}) - 1$.
\end{definition}

\noindent\textbf{Axiom A1 (Normalization):} $J(1) = 0$.\\
\textbf{Axiom A2 (Recognition Composition Law, RCL):}
$J(xy) + J(x/y) = 2J(x)J(y) + 2J(x) + 2J(y)$ for all $x, y > 0$.\\
\textbf{Axiom A3 (Calibration):} $J''_{\log}(0) = 1$, where
$J_{\log}(t) := J(e^t)$.

\begin{theorem}[Cost Uniqueness]
\label{thm:unique}
The unique continuous solution to \textup{(A1)--(A3)} is
$J(x) = \tfrac{1}{2}(x + x^{-1}) - 1$.
\end{theorem}

\begin{proof}[Proof sketch]
Setting $f(t) = J_{\log}(t) + 1$ and rewriting \textup{(A2)} gives the
d'Alembert equation $f(s+t) + f(s-t) = 2f(s)f(t)$.  By the Acz\'el
classification theorem~\cite{Aczel1966}, the unique continuous solution
with $f(0) = 1$ and $f''(0) = 1$ is $f(t) = \cosh t$.
\end{proof}

\begin{theorem}[$\varphi$ Forcing]
\label{thm:phi}
The golden ratio $\varphi = (1+\sqrt{5})/2$ is the unique positive real
satisfying the self-similarity equation from the RCL\@.  Stable masses
lie on the $\varphi$-ladder: $m(r) = \mathrm{yardstick}\cdot\varphi^{r-8+\mathrm{gap}(Q)}$,
where $r$ is the rung index, $8 = 2^D$ for $D=3$, and the gap function
encodes charge dependence.
\end{theorem}

\noindent The $D = 3$ cube has combinatorial integers $V = 8$, $E = 12$, $F = 6$,
$E_{\rm passive} = 11$, and $W = 17$ (wallpaper groups).  Generation spacings
are $\tau(1) = E_{\rm passive} = 11$ and $\tau(2) = W = 17$, with
the inter-generation step gen2$\to$gen3 equal to $\tau(2)-\tau(1) = 6$
(the number of faces of the cube).

%% ============================================================
\section{Rung Assignment}
\label{sec:rung}
%% ============================================================

\begin{definition}[Generation spacings]
\label{def:tau}
The generation spacings $\tau(1)$ and $\tau(2)$ are forced by the
$D = 3$ cube geometry:
\begin{itemize}[nosep]
  \item $\tau(1) = 11$: the number of passive edges in the $D = 3$
  cube (12 total edges minus 1 active edge).
  \item $\tau(2) = 17$: the number of distinct wallpaper groups of
  the $D = 3$ cube.
\end{itemize}
\end{definition}

\begin{proposition}[Top Quark Rung --- Structural Assignment]
\label{thm:rung}
The top quark sits at rung
\begin{equation}
  r_t = 4 + \tau(2) = 4 + 17 = 21,
\end{equation}
where $4$ is the base up-type rung and $\tau(2) = 17$ is the second
generation spacing (wallpaper groups of the $D = 3$ cube).
\end{proposition}

\begin{proof}
The up-type sector has base rung $4$.  First generation: $r_u = 4$.
Second: $r_c = 4 + \tau(1) = 15$.  Third: $r_t = 4 + \tau(2) = 21$.
\end{proof}

\begin{corollary}[Quark rung structure]
The up-type and down-type quarks share the same rung pattern within
each generation: $r_u = r_d = 4$, $r_c = r_s = 15$, $r_t = r_b = 21$.
\end{corollary}

%% ============================================================
\section{Mass Formula}
\label{sec:mass}
%% ============================================================

\begin{proposition}[Top Quark Mass --- RS Prediction]
\label{thm:mass}
\begin{equation}
  m_t = \mathrm{yardstick}(\mathrm{UpQuark}) \cdot
  \varphi^{r_t - 8 + \mathrm{gap}(Q_t)},
\end{equation}
where $\mathrm{yardstick}(\mathrm{UpQuark}) = 2^{-1}
E_{\mathrm{coh}} \varphi^{35}$ and
$\mathrm{gap}(Q_t) = \ln(1 + Q_t/\varphi)/\ln\varphi$ with
$Q_t = 2/3$.
\end{proposition}

\begin{remark}[Gap function conventions]
The gap function used here,
$\mathrm{gap}(Q) = \ln(1 + Q/\varphi)/\ln\varphi$, is a continuous
function of the charge $Q$.  The companion paper on Three Generations
uses the discrete form $\mathrm{gap}(Z) = Z + Z^2 + Z^4$ with
$Z = 6Q$.  Both encode the charge dependence of the RS mass law;
the difference reflects the level of approximation (exact continuous
vs.\ discrete polynomial).  The numerical predictions are similar but
not identical; unifying the two conventions is an identified open problem.
\end{remark}

The exponent $r_t - 8 + \mathrm{gap}(Q_t)$ combines the rung position,
the canonical offset $8$ (from the eight-tick epoch), and the
charge-dependent gap.  The experimental value $m_t = 172.69~\mathrm{GeV}$
is recovered through the calibration seam that fixes the absolute
scale.

%% ============================================================
\section{Quark Mass Hierarchy}
\label{sec:hierarchy}
%% ============================================================

Table~\ref{tab:hierarchy} summarizes the quark mass hierarchy.  RS
assigns integer rungs from cube geometry; the experimental masses
follow the $\varphi$-ladder structure with percent-level deviations
attributed to radiative corrections (QCD running, electromagnetic
effects).

\begin{table}[ht]
\centering
\renewcommand{\arraystretch}{1.3}
\begin{tabular}{lcccc}
\toprule
\textbf{Quark} & \textbf{Rung} & \textbf{Generation} &
\textbf{RS prediction} & \textbf{Exp.\ mass} \\
\midrule
$u$ & 4 & 1st & $\varphi^{-4 + \mathrm{gap}(2/3)}$ & $2.2~\mathrm{MeV}$ \\
$c$ & 15 & 2nd & $\varphi^{7 + \mathrm{gap}(2/3)}$ & $1270~\mathrm{MeV}$ \\
$t$ & 21 & 3rd & $\varphi^{13 + \mathrm{gap}(2/3)}$ & $172690~\mathrm{MeV}$ \\
\midrule
$d$ & 4 & 1st & $\varphi^{-4 + \mathrm{gap}(-1/3)}$ & $4.7~\mathrm{MeV}$ \\
$s$ & 15 & 2nd & $\varphi^{7 + \mathrm{gap}(-1/3)}$ & $93~\mathrm{MeV}$ \\
$b$ & 21 & 3rd & $\varphi^{13 + \mathrm{gap}(-1/3)}$ & $4180~\mathrm{MeV}$ \\
\bottomrule
\end{tabular}
\caption{Quark mass hierarchy: RS rung assignments and experimental
values from PDG~\cite{PDG}.}
\label{tab:hierarchy}
\end{table}

\begin{remark}
The mass ratios between generations (e.g., $m_c/m_u$, $m_t/m_c$)
scale as powers of $\varphi$ determined by the rung differences
$\tau(1) = 11$ and $\tau(2) - \tau(1) = 6$.  The hierarchy is
geometric, not fine-tuned.
\end{remark}

%% ============================================================
\section{Why $y_t \approx 1$: The Top's Natural Scale}
\label{sec:yt}
%% ============================================================

In the Standard Model, $m_t = y_t v/\sqrt{2}$, so $y_t \approx 1$
implies $m_t \sim v$.  RS provides a structural explanation: the top
quark sits at rung 21, which places it near the electroweak scale on
the $\varphi$-ladder.  The Higgs VEV $v \approx 246~\mathrm{GeV}$
occupies a rung $r_v$ such that $v/m_e \approx \varphi^{27}$; the top
mass $m_t \approx 173~\mathrm{GeV}$ is only a few rungs below $v$.
Thus $m_t/v \sim 0.7$ is not a fine-tuning---it is the geometric
consequence of the top occupying the highest fermion rung in the
up-type sector, one step below the scale at which electroweak
symmetry breaks.

\begin{remark}[No hierarchy problem for the top]
In RS, $y_t \approx 1$ is derived: the top's rung is fixed by
$\tau(2) = 17$, and the resulting mass falls naturally near $v$
because both scales sit on the same $\varphi$-ladder.
\end{remark}

%% ============================================================
\section{Top-Bottom Mass Ratio}
\label{sec:tb}
%% ============================================================

The top and bottom quarks occupy the same generation (third) and thus
the same rung $r_t = r_b = 21$.  Their mass difference arises from
the different sector yardsticks and charge gaps.  The ratio
$m_t/m_b$ is predicted by RS to be approximately $\varphi^8$:
\begin{equation}
  \frac{m_t}{m_b} \approx \varphi^8 \approx 46.98.
\end{equation}
Experimentally, $m_t/m_b \approx 41.3$.  The $\sim 12\%$ deviation
matches percent-level corrections in other RS ratios (QCD running,
scheme dependence).

\begin{corollary}
The third-generation up-down mass ratio $m_t/m_b$ is a pure
$\varphi$-power (up to corrections), fixed by the sector structure
without additional parameters.
\end{corollary}

%% ============================================================
\section{Predictions and Falsifiers}
\label{sec:predictions}
%% ============================================================

\begin{enumerate}
  \item \textbf{No fourth generation}: The $D = 3$ cube forces
  exactly three generations.  A confirmed fourth-generation up-type
  quark would falsify the rung derivation.

  \item \textbf{Integer rung}: The top occupies rung 21 exactly.
  The generation spacing $\tau(2) = 17$ is structural, not
  approximate.

  \item \textbf{Top-bottom ratio}: $m_t/m_b$ should lie in the range
  $\varphi^7$ to $\varphi^9$ ($\sim 29$--$76$).  A ratio outside this
  range (after accounting for scheme and scale) would challenge the
  sector structure.

  \item \textbf{Yukawa is derived}: $y_t$ is not a free parameter;
  it is determined by the $\varphi$-ladder rung assignment and the
  calibration seam.

  \item \textbf{Falsifier}: A top mass inconsistent with rung 21
  (e.g., $m_t < 100~\mathrm{GeV}$ or $m_t > 250~\mathrm{GeV}$ at
  the relevant scale) would falsify the prediction.
\end{enumerate}

%% ============================================================
\section{Conclusion}
\label{sec:conclusion}
%% ============================================================

The top quark mass is determined by $\varphi$-rung 21 with no free
Yukawa coupling.  The top's special status---$y_t \approx 1$, mass
near the electroweak scale---is a structural consequence of its
position on the ladder.  The generation spacings $\tau(1) = 11$ and
$\tau(2) = 17$ from $D = 3$ cube geometry fix the entire quark mass
hierarchy.

%% ============================================================
\begin{thebibliography}{99}

\bibitem{Aczel1966}
J.~Acz\'el, \emph{Lectures on Functional Equations and Their Applications},
Academic Press, New York (1966).

\bibitem{RS_Axioms}
J.~Washburn, ``The Algebra of Reality: A Recognition Science Derivation
of Physical Law,'' \emph{Axioms} \textbf{15}(2), 90 (2026).
\texttt{doi:10.3390/axioms15020090}

\bibitem{CDF1995}
F.~Abe \textit{et al.} (CDF Collaboration), Phys.\ Rev.\ Lett.\
\textbf{74}, 2626 (1995).

\bibitem{D01995}
S.~Abachi \textit{et al.} (D0 Collaboration), Phys.\ Rev.\ Lett.\
\textbf{74}, 2632 (1995).

\bibitem{PDG}
R.~L.~Workman \textit{et al.} (Particle Data Group), Prog.\ Theor.\
Exp.\ Phys.\ \textbf{2022}, 083C01 (2022).

\end{thebibliography}

\end{document}
